% ============================================================
% Shared preamble for the RSA module's lecture files
% (lectureN-rsa-*.tex and rsa-course-notes.tex). \input this once
% per document.
%
% Defines: packages and the theorem-style environments used
% throughout the lecture bodies (theorem, lemma, corollary,
% proposition, definition, example, remark, exercise). Deliberately
% plain formatting -- no decorative call-out boxes -- so the
% lectures read as clean, standard mathematical exposition.
%
% Numbering: theorem-family environments are numbered "within
% section" ([section]), so they read as e.g. "Theorem 2.3" inside
% a standalone lectureN.tex (plain \section counter) and as e.g.
% "Theorem 2.2.3" inside rsa-course-notes.tex (book class nests
% \thesection under \thechapter automatically) -- no manual
% renumbering needed between the two contexts.
% ============================================================
\usepackage[margin=1in]{geometry}
\usepackage{amsmath,amssymb,amsthm}
\usepackage{algorithm}
\usepackage{algpseudocode}
\usepackage{enumitem}
\usepackage{hyperref}
\usepackage{booktabs}
\usepackage{xcolor}
\usepackage{titlesec}

% Single shared definition of the RSA module's series title. Every
% document in this project (each standalone lectureN-*.tex wrapper,
% and the combined rsa-course-notes.tex) uses \seriestitle inside its
% own \title{...} command instead of retyping the words -- so the
% title text itself is defined in exactly ONE place. Change it here
% once and every document picks it up on its next compile.
\newcommand{\seriestitle}{The RSA Cryptosystem: Foundations, Attacks, and Implementation}

% Lets \ref (not hardcoded numbers) resolve section/theorem labels that
% live in a DIFFERENT lecture's .tex file -- e.g. Lesson 2 citing a
% theorem number from Lesson 1. Each lecture wrapper file declares which
% other lecture(s) it needs via \externaldocument{<other-lecture-basename>}
% right after % ============================================================
% Shared preamble for the RSA module's lecture files
% (lectureN-rsa-*.tex and rsa-course-notes.tex). \input this once
% per document.
%
% Defines: packages and the theorem-style environments used
% throughout the lecture bodies (theorem, lemma, corollary,
% proposition, definition, example, remark, exercise). Deliberately
% plain formatting -- no decorative call-out boxes -- so the
% lectures read as clean, standard mathematical exposition.
%
% Numbering: theorem-family environments are numbered "within
% section" ([section]), so they read as e.g. "Theorem 2.3" inside
% a standalone lectureN.tex (plain \section counter) and as e.g.
% "Theorem 2.2.3" inside rsa-course-notes.tex (book class nests
% \thesection under \thechapter automatically) -- no manual
% renumbering needed between the two contexts.
% ============================================================
\usepackage[margin=1in]{geometry}
\usepackage{amsmath,amssymb,amsthm}
\usepackage{algorithm}
\usepackage{algpseudocode}
\usepackage{enumitem}
\usepackage{hyperref}
\usepackage{booktabs}
\usepackage{xcolor}
\usepackage{titlesec}

% Single shared definition of the RSA module's series title. Every
% document in this project (each standalone lectureN-*.tex wrapper,
% and the combined rsa-course-notes.tex) uses \seriestitle inside its
% own \title{...} command instead of retyping the words -- so the
% title text itself is defined in exactly ONE place. Change it here
% once and every document picks it up on its next compile.
\newcommand{\seriestitle}{The RSA Cryptosystem: Foundations, Attacks, and Implementation}

% Lets \ref (not hardcoded numbers) resolve section/theorem labels that
% live in a DIFFERENT lecture's .tex file -- e.g. Lesson 2 citing a
% theorem number from Lesson 1. Each lecture wrapper file declares which
% other lecture(s) it needs via \externaldocument{<other-lecture-basename>}
% right after % ============================================================
% Shared preamble for the RSA module's lecture files
% (lectureN-rsa-*.tex and rsa-course-notes.tex). \input this once
% per document.
%
% Defines: packages and the theorem-style environments used
% throughout the lecture bodies (theorem, lemma, corollary,
% proposition, definition, example, remark, exercise). Deliberately
% plain formatting -- no decorative call-out boxes -- so the
% lectures read as clean, standard mathematical exposition.
%
% Numbering: theorem-family environments are numbered "within
% section" ([section]), so they read as e.g. "Theorem 2.3" inside
% a standalone lectureN.tex (plain \section counter) and as e.g.
% "Theorem 2.2.3" inside rsa-course-notes.tex (book class nests
% \thesection under \thechapter automatically) -- no manual
% renumbering needed between the two contexts.
% ============================================================
\usepackage[margin=1in]{geometry}
\usepackage{amsmath,amssymb,amsthm}
\usepackage{algorithm}
\usepackage{algpseudocode}
\usepackage{enumitem}
\usepackage{hyperref}
\usepackage{booktabs}
\usepackage{xcolor}
\usepackage{titlesec}

% Single shared definition of the RSA module's series title. Every
% document in this project (each standalone lectureN-*.tex wrapper,
% and the combined rsa-course-notes.tex) uses \seriestitle inside its
% own \title{...} command instead of retyping the words -- so the
% title text itself is defined in exactly ONE place. Change it here
% once and every document picks it up on its next compile.
\newcommand{\seriestitle}{The RSA Cryptosystem: Foundations, Attacks, and Implementation}

% Lets \ref (not hardcoded numbers) resolve section/theorem labels that
% live in a DIFFERENT lecture's .tex file -- e.g. Lesson 2 citing a
% theorem number from Lesson 1. Each lecture wrapper file declares which
% other lecture(s) it needs via \externaldocument{<other-lecture-basename>}
% right after % ============================================================
% Shared preamble for the RSA module's lecture files
% (lectureN-rsa-*.tex and rsa-course-notes.tex). \input this once
% per document.
%
% Defines: packages and the theorem-style environments used
% throughout the lecture bodies (theorem, lemma, corollary,
% proposition, definition, example, remark, exercise). Deliberately
% plain formatting -- no decorative call-out boxes -- so the
% lectures read as clean, standard mathematical exposition.
%
% Numbering: theorem-family environments are numbered "within
% section" ([section]), so they read as e.g. "Theorem 2.3" inside
% a standalone lectureN.tex (plain \section counter) and as e.g.
% "Theorem 2.2.3" inside rsa-course-notes.tex (book class nests
% \thesection under \thechapter automatically) -- no manual
% renumbering needed between the two contexts.
% ============================================================
\usepackage[margin=1in]{geometry}
\usepackage{amsmath,amssymb,amsthm}
\usepackage{algorithm}
\usepackage{algpseudocode}
\usepackage{enumitem}
\usepackage{hyperref}
\usepackage{booktabs}
\usepackage{xcolor}
\usepackage{titlesec}

% Single shared definition of the RSA module's series title. Every
% document in this project (each standalone lectureN-*.tex wrapper,
% and the combined rsa-course-notes.tex) uses \seriestitle inside its
% own \title{...} command instead of retyping the words -- so the
% title text itself is defined in exactly ONE place. Change it here
% once and every document picks it up on its next compile.
\newcommand{\seriestitle}{The RSA Cryptosystem: Foundations, Attacks, and Implementation}

% Lets \ref (not hardcoded numbers) resolve section/theorem labels that
% live in a DIFFERENT lecture's .tex file -- e.g. Lesson 2 citing a
% theorem number from Lesson 1. Each lecture wrapper file declares which
% other lecture(s) it needs via \externaldocument{<other-lecture-basename>}
% right after \input{lecture-preamble}. This requires the referenced
% lecture to have been compiled at least once already (so its .aux
% exists) -- compile_all.sh handles this automatically by building the
% lectures in order (1, then 2, then 3). If you compile a single later
% lecture in isolation before ever compiling the earlier one(s) it cites,
% those specific cross-references will show as "??" until you compile
% the earlier lecture(s) once.
\usepackage{xr-hyper}

\newtheorem{theorem}{Theorem}[section]
\newtheorem{lemma}[theorem]{Lemma}
\newtheorem{corollary}[theorem]{Corollary}
\newtheorem{proposition}[theorem]{Proposition}
\newtheorem{definition}[theorem]{Definition}
\newtheorem{example}[theorem]{Example}
\newtheorem{remark}[theorem]{Remark}
\newtheorem{exercise}[theorem]{Exercise}

\hypersetup{
    colorlinks=true,
    linkcolor=blue!50!black,
    citecolor=blue!50!black,
    urlcolor=blue!50!black,
    pdfpagemode=UseNone
}
. This requires the referenced
% lecture to have been compiled at least once already (so its .aux
% exists) -- compile_all.sh handles this automatically by building the
% lectures in order (1, then 2, then 3). If you compile a single later
% lecture in isolation before ever compiling the earlier one(s) it cites,
% those specific cross-references will show as "??" until you compile
% the earlier lecture(s) once.
\usepackage{xr-hyper}

\newtheorem{theorem}{Theorem}[section]
\newtheorem{lemma}[theorem]{Lemma}
\newtheorem{corollary}[theorem]{Corollary}
\newtheorem{proposition}[theorem]{Proposition}
\newtheorem{definition}[theorem]{Definition}
\newtheorem{example}[theorem]{Example}
\newtheorem{remark}[theorem]{Remark}
\newtheorem{exercise}[theorem]{Exercise}

\hypersetup{
    colorlinks=true,
    linkcolor=blue!50!black,
    citecolor=blue!50!black,
    urlcolor=blue!50!black,
    pdfpagemode=UseNone
}
. This requires the referenced
% lecture to have been compiled at least once already (so its .aux
% exists) -- compile_all.sh handles this automatically by building the
% lectures in order (1, then 2, then 3). If you compile a single later
% lecture in isolation before ever compiling the earlier one(s) it cites,
% those specific cross-references will show as "??" until you compile
% the earlier lecture(s) once.
\usepackage{xr-hyper}

\newtheorem{theorem}{Theorem}[section]
\newtheorem{lemma}[theorem]{Lemma}
\newtheorem{corollary}[theorem]{Corollary}
\newtheorem{proposition}[theorem]{Proposition}
\newtheorem{definition}[theorem]{Definition}
\newtheorem{example}[theorem]{Example}
\newtheorem{remark}[theorem]{Remark}
\newtheorem{exercise}[theorem]{Exercise}

\hypersetup{
    colorlinks=true,
    linkcolor=blue!50!black,
    citecolor=blue!50!black,
    urlcolor=blue!50!black,
    pdfpagemode=UseNone
}
. This requires the referenced
% lecture to have been compiled at least once already (so its .aux
% exists) -- compile_all.sh handles this automatically by building the
% lectures in order (1, then 2, then 3). If you compile a single later
% lecture in isolation before ever compiling the earlier one(s) it cites,
% those specific cross-references will show as "??" until you compile
% the earlier lecture(s) once.
\usepackage{xr-hyper}

\newtheorem{theorem}{Theorem}[section]
\newtheorem{lemma}[theorem]{Lemma}
\newtheorem{corollary}[theorem]{Corollary}
\newtheorem{proposition}[theorem]{Proposition}
\newtheorem{definition}[theorem]{Definition}
\newtheorem{example}[theorem]{Example}
\newtheorem{remark}[theorem]{Remark}
\newtheorem{exercise}[theorem]{Exercise}

\hypersetup{
    colorlinks=true,
    linkcolor=blue!50!black,
    citecolor=blue!50!black,
    urlcolor=blue!50!black,
    pdfpagemode=UseNone
}
